\section{Introduction}
\label{sec:introduction}

Differential Fault Analysis recovers cryptographic secrets by inducing
transient faults during cipher execution and comparing corrupted outputs
against correct ones~\cite{biham1997dfa,bareltal2006fault}. A natural
way to characterize a hardware implementation is therefore to inject a
single fault at each gate output and measure how often that fault reaches
a primary output. Recent synthesis-aware fault studies use exactly this
kind of gate-level campaign to compare representations such as LUT
netlists, AND-inverter graphs, majority-inverter graphs, and
technology-mapped FPGA circuits.

This paper asks a simple methodological question: when such a campaign
reports that one synthesis representation has lower per-vector fault
observability than another, how much of the effect is caused by the
logic primitive, and how much is caused by the granularity of the
netlist being measured? The distinction matters. A wide truth-table LUT
that directly drives a primary output contributes one observable fault
site with detection probability one. Decomposing the same function into
many small gates creates many interior sites whose effects can be masked
downstream. A metric averaged over all gate sites can therefore move
substantially even when the realized Boolean function and attacker-facing
target layer have not become more fault resistant.

We evaluate this question on 15 cipher primitives synthesized into three
representations: the LUT-style AOIG baseline produced by Yosys, a
decomposed 2-input AIG as a comparable small-gate control, and a
3-input MIG. All three are subjected to the same single gate-output
\bitflip{} campaign. The resulting observation is narrow but important:
the reported AOIG-to-MIG \AvgDP{} reduction is reproducible, but it
does not survive a decomposed-AIG baseline at comparable granularity.

\subsection*{Contributions}

\begin{itemize}
  \item We show that the AOIG-to-MIG \AvgDP{} reduction reverses sign
        under a decomposed-AIG baseline at comparable small-gate granularity:
        LUT$\to$MIG gives a mean $13.0\%$ reduction, while AIG$\to$MIG
        gives a mean $21.4\%$ increase with $0/15$ circuits reduced.
  \item As a granularity control, we measure target-layer \AvgDP{}
        (faults restricted to output-driving gates). This yields \AvgDP{}
        $=1.000$ across all three representations, consistent with the
        interpretation that the all-site difference comes from interior
        fault-site population mix, not from any change in the
        attacker-facing output layer.
  \item We report one result that survives the controls: ARX primitives exhibit
        approximately $1.9\times$ broader output corruption per detected
        fault than substitution-based primitives across representations.
  \item We release the full artifact: netlists, fault simulator, frozen
        JSON outputs, generated tables/figures, and a script that checks
        the paper's quoted summary numbers.\footnote{Artifact:
        \url{https://github.com/Mrunal321/cipher-fault-observability-study}}
\end{itemize}

The rest of the paper proceeds as follows. \S\ref{sec:background}
summarizes MIG synthesis and DFA observability metrics.
\S\ref{sec:methodology} defines the benchmark, fault model, and
granularity-controlled methodology. \S\ref{sec:results} gives the
baseline-flip result, the invariant target-layer result, and the stable
architecture-level finding. \S\ref{sec:discussion} summarizes
reporting guidance and remaining limits, and \S\ref{sec:conclusion}
concludes.
